\documentclass[]{article}

\usepackage[a5paper]{geometry}
\usepackage[T1]{fontenc}
\usepackage[italian]{babel}
\usepackage{tikz}


\begin{document}

\title{\Huge \bfseries Truss optimizer}
\author{Andrea Marchi}
\maketitle


\section{Ottimizzazione aree elementi}

L'idea e` quella che iterativamente il FEM trova la sezione necessaria $A$ per ottenere la tensione obiettivo $\sigma_{ob}$.
Risolto per la prima volta il sistema, l'elemento $i$ e` sottoposto alla tensione
$$ \sigma_i(t_1) = \frac{N_i(t_1)}{A_i(t_1)} $$
dove $t_1$ indica che stiamo considerando il primo istante, l'istante $1$.
A questo punto dobbiamo modificare l'area degli elementi $A_i$ per avvicinarci alla tensione obiettivo
$$ A_i(t_2) = \frac{|N_i(t_1)|}{\sigma_{ob}} \geq A_{\min} $$
Per evitare che in aste completamente scariche si ha un'area nulla bisogna comunque avere un vincolo di area minima $A_{\min}$.
Una volta settate queste nuove aree il sistema viene nuovamente risolto per trovare i nuovi sforzi nelle aste. Questo, in generale, porta a sforzi diversi e quindi aree diverse per le iterazioni successive.
Le iterazioni sono portate avanti finche` il sistema non arriva a convergenza. Questo puo` misurarsi in termini di differenze della somma delle $\sigma$ al quadrato per tutti gli elementi e vedendo la differenza tra una iterazione e l'altra. Quindi il programma termina il calcolo quando
$$ \sum_{i=1}^N \left[|\sigma_i(t_j)| - \sigma_{ob} \right]^2 \leq \mathrm{Tol} $$

Ovviamente i sistemi isostatici non hanno bisogno di iterazioni e le aree dopo la prima iterazione sono gia` quelle definitive.

\textcolor{red}{Se c'e' un elemento scarico questa metodologia non converge...}
Questo perche' $\sigma_i = N_i / A_i$, dato che $N_i$ e' pari a zero, allora si somma all'errore un termine $N_{zero} \; \sigma_{ob}^2$ ($N_{zero}$ e' il numero di elementi in cui $N=0$).
Per ovviare a questo potrei definire l'errore in termini di differenza percentuale tra aree (piu' complesso da programmare)
$$ \mathrm{errore}_j := \sum_{i=1}^N \frac{|A_i(t_{j+1}) - A_i(t_j)|}{A_i(t_{j+1})} $$
dove $\mathrm{errore}_j$ e' l'errore all'istante $j$ e la somma e' fatta su tutti gli $N$ elementi.
Ovviamente l'errore deve essere minore di una tolleranza, considerata in termini percentuali rispetto all'area.

La velocita' di convergenza sembra essere lineare e dipende dal numero di iperstatiche (maggiore iperstaticita' e maggior tempo ci vuole per far convergere il risultato, ovvero la velocita' e' minore).


\section{Ottimizzazione struttura}
\subsection{Eliminazione elementi}

In questa strategia partiamo da una struttura a molte iperstatiche e, ad ogni iterazione, eliminiamo un elemento, fino a raggiungere una struttura ottimizzata.

\textcolor{red}{Per capire se ha eliminato tanti elementi da ottenere un sistema isostatico potrei vedere il numero di iterazioni necessarie per ottimizzare le aree. Se il numero e' pari a 1 allora ho raggiunto una struttura isostatica...}

Bisogna eliminare i nodi e gli elementi "degeneri":
\begin{itemize}
\item nodi degeneri: quei nodi che non sono collegati a nessun elemento.
\item elementi degeneri: aste che "penzolano". Sono elementi collegati a nodi che presentano un solo elemento e niente vincoli.
\end{itemize}





\end{document}